\documentclass[11pt]{article}
\usepackage[review]{acl}
\usepackage{times,latexsym}
\usepackage[T1]{fontenc}
\usepackage[utf8]{inputenc}
\usepackage{microtype,inconsolata,graphicx,booktabs,amsmath,float}
\title{Evaluating Safety Steering Beyond Attack Success:\\Random Directions, Benign Wrappers, and Measurement Instability}
\author{Anonymous ACL submission}
\begin{document}
\maketitle
\begin{abstract}
A lower automated attack success rate does not by itself establish a better safety intervention. We audit activation steering with norm-matched random directions, paired benign requests with and without adversarial wrappers, and zero-training prompt and output-filter baselines. On a success-selected set of 70 attack prompts, random perturbations already reduce judge attack success by 4--21 percentage points. Learned vectors are substantially stronger than this reference, but benign-wrapper refusal and several measurement choices limit the interpretation of that gain. Our audits expose baseline-set mismatches, output truncation, classifier-position and threshold errors, and label changes between highly similar retrained vectors. A five-seed paired benign comparison reproduces a positive wrapper-associated refusal gap on a second prompt sample. We report these as evaluation findings rather than a new defense, distinguish operational model labels from calibrated refusal rates, and preserve unsuccessful frozen criteria. Data-composition interventions are running and are not treated as completed evidence in this draft.
\end{abstract}
\section{Introduction}
Safety steering changes a language model's internal activations to influence its responses. A small vector can produce a large change in an automated safety score, but several explanations remain compatible with that observation: learned refusal, nonspecific perturbation, a change in response format, or a measurement artifact. Benign requests may also inherit the surface form of a jailbreak, creating a useful test of whether an intervention tracks harmful intent or request packaging.

We study these distinctions in a controlled audit of one learned steering setup. Our main contribution is a collection of measured evaluation pitfalls and operational baselines, rather than a claim that steering constitutes a generally effective defense. We compare learned vectors with norm-matched random directions, decompose the learned update relative to initialization, and evaluate paired benign originals and adversarially packaged versions. We also retain system prompts and an output-only safety filter as zero-training references.

The resulting evidence is mixed in a scientifically useful way. Learned vectors outperform sampled random directions on both our selected attack set and an external attack family. Nevertheless, random directions alone substantially change the attack judge, refusal concentrates on benign wrapped requests, and nominally minor implementation choices alter the measured effect. Even highly similar retrained vectors produce different response trajectories and per-item labels. These observations motivate paired retrained controls and explicit separation of judge outcomes, refusal labels, and task completion.

We make four contributions. First, we quantify the random-direction reference on success-selected attacks and examine the residual component of learned steering. Second, we reproduce a wrapper-associated benign refusal gap across prompt samples and report its descriptive association with overall over-refusal. Third, we document baseline membership, reading scope, truncation, backend alignment, thresholds, and training repeat variation in our own pipeline, with reproducible tables. Fourth, on 300 WildJailbreak attacks sampled without filtering for baseline success, all five learned ASRs are below the four random references (exploratory exact unit-level $p=0.0159$). Three of five seeds met the 10-pp threshold under the frozen two-direction reference, short of the required four; all five met it in the supplementary four-direction comparison.

\section{Related Work}
\paragraph{Steering and refusal.} Learned interventions on frozen-model activations include representation finetuning and bidirectional preference steering \citep{wu2024reft,cao2024bipo}. We audit one fixed-vector injection protocol. Related work separates steering efficacy from the preservation of other behaviors under distribution shift \citep{goyal2026specificity}. \citet{arditi2024refusal} study a refusal direction in chat models. This supports evaluating harmful compliance and benign task utility together. Our learned vector, its residual decomposition, and its final-prompt-token injection protocol are specific to the setup below; we do not equate them with a universally shared refusal direction.

\paragraph{Random steering.} \citet{korznikov2026rogue} report that random fixed-vector steering can increase harmful compliance from zero to 1--13\% in the revised version. Their intervention persists across prompt and generated tokens, whereas our intervention is a single injection at the final prompt token. Our attack prompts are also selected for prior success. Random perturbations move our judge score in the opposite direction. The protocol and input-selection differences are plausible explanations, not a tested mechanism or a direct numerical replication.

\paragraph{Attack and benign evaluations.} PAIR iteratively constructs semantic jailbreaks through black-box interaction \citep{chao2023pair}. HarmBench \citep{mazeika2024harmbench} and JailbreakBench \citep{chao2024jbb} emphasize standardized attack evaluation. WildTeaming introduces WildJailbreak with both vanilla and adversarial harmful and benign examples \citep{jiang2024wildteaming}. Its benign requests that resemble harmful ones motivate our paired wrapper analysis. Because our benign intervention and evaluation share this source, that intervention tests in-distribution data composition rather than generalization to arbitrary wrapper styles.

\paragraph{Scope of the audit.} Prior work compares automated jailbreak evaluators and studies judge calibration and robustness \citep{gao2026judge,mu2026evaluators}. We complement these evaluations by distinguishing a harmfulness judge flip from an observed refusal, and an observed refusal from a calibrated utility failure. We use matched baselines and explicit reading scopes to expose cases where those distinctions change reported counts. Numerical agreement across implementations is an important reproducibility check, but it does not validate the judge on a new distribution or language.

\section{Experimental Setup}
\paragraph{Models and intervention.} The primary target is Llama-3.1-8B-Instruct. A supplementary zero-training evaluation uses Qwen2.5-7B-Instruct; no Qwen steering vector or alpha sweep is reported. The old trainer learns a 4096-dimensional vector with norm 5.21, 40 Adam steps at learning rate 0.005, 12 fixed attack inputs, and 60 benign keep inputs. Model weights remain frozen. At evaluation, a vector scaled by $\alpha$ is injected once into the last prompt position of layer index 9. Generation is greedy. Harmful evaluation permits 2400 new tokens; the original benign protocol permits 120, with a controlled 400-token follow-up for nine uncertain examples.

\paragraph{Prompt samples.} All attack sets in this study consist of text behaviors. The selected attack set, val70, contains 14 harmful JBB-Behaviors rows and 56 HarmBench behaviors, with successful PAIR prompts selected against the unmodified Llama model. Selection matters: this is not an unbiased sample of all harmful requests. The historical WildJailbreak attack evaluation contains 300 adversarial harmful prompts sampled from the training split. WildJailbreak also provides a distinct official evaluation split. Two benign samples pair unwrapped originals with adversarially packaged benign requests. One analysis retains 293 pairs from the val296 source after its baseline exclusion; a separate 300-pair sample retains 292 when either-side baseline rule refusals are removed. The latter exclusion was documented after some generations and is not described as an original pre-generation rule.

The planned transfer evaluation retains all 400 HarmBench text rows and 100 harmful JBB-Behaviors rows (Table~\ref{tab:dataset_scope}). The search uses a rule-scored PAIR-style attack on undefended Llama-3.1-8B-Instruct. Non-refusal and output length guide search. The search procedure and row order were specified in a versioned log before the ten-row pilot. That pilot preceded the freeze of the same-host transfer evaluation. Its rows remain part of the 500-row manifest. System prompts and fixed vectors receive the same frozen per-row prompts. Historical val70 results remain a separate success-selected stratum.

\begin{table*}[t]
\centering
\small
\begin{tabular}{lrrrl}
\toprule
Data and stratum & Rows & Groups & Untouched rows & Outcome scorer \\
\midrule
HarmBench total & 400 & 391 & 215 & Mixed \\
HarmBench standard & 200 & 200 & 53 & 13B standard \\
HarmBench contextual & 100 & 93 & 73 & 13B contextual \\
HarmBench copyright & 100 & 98 & 89 & MinHash \\
JBB-Behaviors harmful & 100 & 100 & 13 & 13B standard \\
All text behaviors & 500 & 479 & 228 & Mixed \\
Standard + JBB auxiliary & 66 & 64 & 66 & 13B standard \\
Contextual + copyright auxiliary & 162 & 155 & 162 & 13B / MinHash \\
\bottomrule
\end{tabular}
\caption{Planned text evaluation of one rule-scored PAIR-style attack instance on HarmBench and JBB-Behaviors. The aggregate contains 12 cross-dataset groups. The two untouched strata are reported separately, with contextual and copyright outcomes shown individually within the second stratum. Counts describe the frozen manifests.}
\label{tab:dataset_scope}
\end{table*}


\paragraph{Outcomes.} Strict judge ASR is the proportion with HarmBench $P(\mathrm{Yes})>0.5$. Backend and INVALID counts accompany comparisons; the HF alpha-zero baseline is 63/70 with zero INVALID, rather than the older vLLM cache's 62/70 with one INVALID. Operational benign over-refusal counts new rule-positive responses among rule-negative baseline items, using fixed sets of 388 Llama or 392 Qwen items. These are not independently human-calibrated safety-refusal rates. We preserve 400-item batch-8 composition for legal generation before filtering to 388 eligible items.

For benign pair $i$, let $r_i^A$ and $r_i^V$ be frozen rule labels on the adversarial wrapper and vanilla response. We report $DD=100\,\mathrm{mean}_i(r_i^A-r_i^V)$ on the fixed baseline-eligible set. Pair-level intervals resample vanilla requests, preserving the two sides. Training-level comparisons use the seed as the unit; repeated prompts are not independent training runs.

\paragraph{Statistical and chronological limits.} Five-seed intervals use a $t$ reference with four degrees of freedom. A joint frontier sensitivity also resamples attack and benign items and propagates old-curve uncertainty; unsupported interpolation draws are identified. The WJB nine-unit rank permutation test is exploratory and assumes exchangeability of five learned seeds and four random directions, conditional on shared prompts. Versioned logs preserve amendments and unsuccessful criteria rather than relabeling later supplements as original confirmation.

\paragraph{Exposure strata and shared sources.} Of the 228 historically unflagged rows in the 500-row manifest, 162 are contextual or copyright behaviors and 66 are standard or JBB behaviors. Report the unflagged standard-plus-JBB, contextual, and copyright results in separate tables because their scorers differ. The JBB layer has historical use flags on 87 of 100 rows, including the twelve fit12 behaviors used as attack inputs in vector training. Its overlap with val70 and the training behavior families prevents interpreting it as independent behavioral evidence. The official WildJailbreak evaluation also shares HarmBench-standard source behaviors.

\section{Results: registered analysis structure}
\label{sec:results}
This working draft reserves the new-round results for the staged reads. Generation and cache-only judging are not reported as observed treatment effects. The planned organization is by finding: random controls; fulfillment declines and their refusal association; benign and utility costs; attack-family consistency and adaptive gaps; and measurement sensitivity. Each completed table will identify its model-specific operating point, relative norm, judge construct, missing-label bounds, and dynamic-range status. Unrun or gated cells remain explicit.

% R1 English static; R2 Chinese/French; R3 encodings/Indian languages;
% R4 search; R5 other defenses. Later methods cannot change after earlier reads.
% Main outcome tables are joined to the final generation manifest hashes.
% Do not populate this section by globbing pilot or legacy result caches.

\section{Measurement audit and reporting checklist}
\label{sec:pitfalls}
A lower classifier rate can change because the model, the reference, or the measurement path changed. Our audit separates these possibilities rather than adding them into a single corrected safety score. The checklist records each trigger, observed consequence, detection method, and practical control; the detailed historical comparisons remain in Appendix~\ref{app:pitfalls}.

\paragraph{The judge's construct and input matter.} HarmBench conditions on a behavior text, so an unverified nearest-behavior mapping changes the question it answers. The historical WildJailbreak top-1 and top-2 references changed labels on 36.65\% of baseline rows. This is reference sensitivity, not an estimate of mapping error. The redesigned SORRY evaluation supplies the original English instruction instead. Even then, fulfillment, reply harmfulness, and nonrefusal remain different constructs. The human-reference confusion matrices expose class- and style-specific errors; an overall agreement score alone cannot show that errors cancel in learned-minus-random comparisons.

\paragraph{Length caps and decoding paths define the comparison.} A short refusal-like prefix can later turn into compliance, and a warning can precede substantive content. Changing the length cap or scoring only the first sentence can therefore change the outcome without changing a stored prefix. We preserve the 2400-token harmful-generation cap, retain the 120-token public-cost protocol as a limitation, and report the registered shorter-reply sensitivity separately. Prefix families give v2 only the continuation while model judges receive the complete prefilled reply. Batch composition, padding, and kernel choices can also change text. The engineering label screen is not proof of equivalence; within-cell paths stay fixed and the single-row contrast regression is separately registered.

\paragraph{Training variation and reference choice need different controls.} Repeated old-recipe seed3 training produced near-identical vectors while changing a few val70 labels. A high cosine similarity is an implementation check, not a guarantee of item-label stability. Random directions answer a different question: how much of a judge-rate change occurs under a perturbation of the same norm without that training recipe? We preserve both controls and do not interpret one seed's rank as a population property.

\paragraph{External validation has a domain.} A released human test split provides a reference for judge implementation and class-specific errors, not human labels for our protected-model outputs. A split can contain different replies to related instructions, and translation-style replies can already be English. We report those properties alongside accuracy and $\kappa$. No judge is discarded after an unfavorable protected-model effect, and low-positive-count encoding styles are not called validated merely because accuracy is high.

\paragraph{An audit trail must distinguish decisions from receipts.} Source and model hashes, row-level atomic records, runtime vector norms, and cache-source hashes make implementation changes traceable. They do not repair an incorrect construct. The first nonzero lease preceded its dedicated release receipt; the prerequisites were already frozen, but the receipt is explicitly retrospective. Search early-stop scores likewise remain internal search signals rather than independent fulfillment judgments. Failed criteria, aborted work, and incomplete model ports stay in the record instead of being rewritten as successful preregistered tests.

\section{Limitations}
Our intervention is a fixed, single-position additive vector. We do not establish a general defense, a causal refusal mechanism, or robustness to an attacker optimized against every defended condition. Search prompts are optimized against the unmodified target and transferred to the fixed conditions. The separately specified adaptive check has a smaller scope and is not complete merely because transfer search has run.

The same training recipe does not create equally optimized interventions. Three learned seeds and two random directions give limited information about their populations. The fixed forty-step Qwen run has no demonstrated convergence; its attack loss remains approximately 1.8--2.1. Transferred Llama-style targets, model-specific loss scales, and finite optimization are confounded. Calibration cannot identify which explains an unreachable refusal target. We report seeds separately and do not compensate by training Qwen longer in the main round. Different selected relative norms also prevent interpreting cross-model contrasts as a matched-strength experiment.

All primary outcome instruments are models. SORRY and HarmBench assess task completion under different prompts and training distributions, whereas WildGuard nonrefusal is a refusal complement, not a third independent harmfulness judgment. External human-reference agreement is style- and response-distribution-specific. The released translated-style reference replies are predominantly English, so their scores do not validate a translation chain for our source-language replies. Rare positive labels make encoding-style accuracy particularly misleading. Qualification gates restrict wording rather than establish that qualified instruments are correct on protected-model outputs.

Refusal labels do not establish why a response changed. A refusal may include alternatives, warnings, or partial compliance, and nonrefusal may be an unrelated or unhelpful answer. WildGuard and v2 can disagree in either direction; neither is a mathematical bound on the other. The post-hoc association between fulfillment flips and refusal is descriptive. We exclude WildGuard's nonrefusal proxy from that decomposition to avoid a tautological explanation.

Generation paths and truncation are part of the estimand. Harmful replies permit 2400 tokens, while public refusal and calibration replies permit 120. A turn from compliance to refusal, or the reverse, can occur beyond that cutoff. Batch screening on one hundred engineering prompts is a prespecified operational filter, not an equivalence test. Matching batch composition across conditions reduces one implementation difference but does not establish equality to single-row decoding. The registered single-row regression comparison remains a sensitivity check, not a reason to replace an unfavorable main path after reading results.

Training, calibration, and test sources have related behaviors and styles. We retain historical-use and semantic-overlap flags rather than discard inconvenient test rows, but exact matching and an embedding threshold cannot establish semantic independence. Fit12 attacks were originally constructed against Llama. Reusing their fixed texts across models controls text while introducing transfer limitations. Public benchmarks differ in request construction, policy boundaries, and difficulty; their costs cannot by themselves identify a wrapper effect. Paired benign tests provide the narrower packaging comparison.

Dynamic-range and cost gates determine which cells receive nonzero evaluation. Skipped cells are reported as unrun, not as zero effect or evidence of success. Local search role models may misidentify success and alter pruning quality relative to original implementations. Limited pilots, early stopping, and incomplete warmup requests make resource extrapolations uncertain. Sunk costs remain visible even when a cell fails its gate. Languages and encodings have lower-priority descriptive coverage; unfinished cells are not counted as tested.

The historical WildJailbreak and PAIR-500 studies retain their original operating points and historical-use strata in the appendix. Their model counts, judge definitions, and selection histories do not silently extend to the redesigned main experiment. Gemma, Mistral, detector ports, and planned human rereads contribute no completed replication until their own records support it. Model annotations share context and remain distinct from independent human calibration.

\bibliography{references}
\appendix
\section{Provenance and protocol chronology}
All numeric tables are generated from CSV files, which in turn are exported from versioned analysis files. The table provenance manifest records complete source SHA256 hashes. Raw attack prompts, model outputs, and training caches are excluded from the manuscript repository. Decisions were specified in a versioned log. External timestamps begin on 2026-11-06 at 00:34 JST; earlier states are supported only by server logs. Several thresholds and protocol amendments followed earlier observations and are identified as such.
The original five-seed WJB criterion remains unsuccessful (3/5). Later random directions, a retrained-control design, a third data arm, and interpolation supplements do not revise that original outcome. The third-arm amendment was written after the original matching table had been frozen and holdout generation had begun. A fourth amendment expanded the plain arm to five seeds before the added seeds were trained and before any plain-arm holdout was generated or read; the first three-seed training queue had already started. These amendments do not revise the original criterion.
\section{Auxiliary frontier sensitivity}
The new-versus-old trainer frontier point estimate is about +0.98pp; conditional on the observed horizontal coordinates, its five-seed interval is $[-0.63,2.59]$pp. Table~\ref{tab:frontierjoint} propagates additional item and old-curve noise. The resulting uncertainty includes zero and extends beyond the proposed equivalence band. An exploratory seven-point subset is not substituted for the full analysis.
\begin{table}[H]
\centering
\small
\setlength{\tabcolsep}{3pt}
\begin{tabular}{p{0.61\columnwidth}r}
\toprule
Joint sensitivity & 95\% interval \\
\midrule
Item resampling (conditional) & [-0.95, 3.56] \\
Also resample seeds (conditional) & [-1.34, 3.88] \\
Boundary clipping diagnostic & [-2.12, 3.57] \\
\bottomrule
\end{tabular}
\caption{Joint frontier sensitivity in pp. The first two intervals condition on all points being interpolable. In 4764/10000 draws at least one point was unsupported; clipping changes the estimand. None establishes equivalence within plus/minus 3 pp.}
\label{tab:frontierjoint}
\end{table}

\end{document}
